% =============================================================================
%  Summary Sheet 正文：直接写段落，段与段之间空行分隔。
%  不要塞回 main.tex 的 \newcommand —— 宏不是 \long，一分段就报
%  「Paragraph ended before \new@command was complete」，且日志不指向摘要。
%
%  五段结构（一页装下）：
%    ① 背景：现象 → 危害 → 本类方案为什么适用（4–6 句，不写方法）
%    ②③④ 每题一段：做法一句 → 关键数值 2–3 组 → 结论一句
%    ⑤ 收尾：模型与策略的整体价值（两句）
%
%  三条铁律：
%    1. 同一指标全文一个量纲；摘要里的每个数都要能在正文找到同一个数。
%    2. 数字方向与结论同向——指标变大就解释为什么，不许写 decreases。
%    3. 抑制率 = (关 − 开)/关，必须能用摘要里的两个数复算出来。
% =============================================================================

[Summary: one paragraph of background; then one paragraph per problem --- the method in
one sentence, two or three numerical results with units, and the conclusion that follows;
then a closing paragraph. 300--500 words, no formulas, no figures, no citations.]
